\documentclass[12pt, a4paper]{article}

% ============================================================
% PACKAGES
% ============================================================
\usepackage[utf8]{inputenc}
\usepackage[T1]{fontenc}
\usepackage{lmodern}
\usepackage[margin=1in]{geometry}
\usepackage{setspace}
\usepackage{graphicx}
\usepackage{booktabs}
\usepackage{longtable}
\usepackage{array}
\usepackage{multirow}
\usepackage{float}
\usepackage{caption}
\usepackage{subcaption}
\usepackage{amsmath, amssymb}
\usepackage{algorithm}
\usepackage{algpseudocode}
\usepackage{listings}
\usepackage{xcolor}
\usepackage{hyperref}
\usepackage{url}
\usepackage{enumitem}
\usepackage{tikz}
\usetikzlibrary{shapes.geometric, arrows.meta, positioning, fit, calc}
\usepackage{pgfplots}
\pgfplotsset{compat=1.18}
\usepackage[backend=biber, style=ieee, sorting=nyt]{biblatex}
\usepackage{glossaries}
\usepackage{appendix}
\usepackage{tcolorbox}
\usepackage{fancyhdr}

% ============================================================
% CONFIGURATION
% ============================================================

% Hyperlink styling
\hypersetup{
    colorlinks=true,
    linkcolor=blue!70!black,
    citecolor=green!50!black,
    urlcolor=blue!70!black,
}

% Code listing style
\definecolor{codebg}{rgb}{0.96, 0.96, 0.96}
\definecolor{codegreen}{rgb}{0, 0.5, 0}
\definecolor{codegray}{rgb}{0.5, 0.5, 0.5}
\definecolor{codepurple}{rgb}{0.58, 0, 0.82}

\lstdefinestyle{pythonstyle}{
    backgroundcolor=\color{codebg},
    basicstyle=\ttfamily\footnotesize,
    breaklines=true,
    captionpos=b,
    commentstyle=\color{codegreen},
    keywordstyle=\color{blue},
    stringstyle=\color{codepurple},
    numberstyle=\tiny\color{codegray},
    numbers=left,
    numbersep=5pt,
    frame=single,
    rulecolor=\color{black!30},
    language=Python,
    showstringspaces=false,
    tabsize=4,
}

\lstdefinestyle{jsonstyle}{
    backgroundcolor=\color{codebg},
    basicstyle=\ttfamily\footnotesize,
    breaklines=true,
    frame=single,
    rulecolor=\color{black!30},
    showstringspaces=false,
}

\lstdefinestyle{sqlstyle}{
    backgroundcolor=\color{codebg},
    basicstyle=\ttfamily\footnotesize,
    breaklines=true,
    captionpos=b,
    commentstyle=\color{codegreen},
    keywordstyle=\color{blue},
    language=SQL,
    frame=single,
    rulecolor=\color{black!30},
    showstringspaces=false,
}

% Header/footer
\pagestyle{fancy}
\fancyhf{}
\fancyhead[L]{\leftmark}
\fancyhead[R]{Ad LLM as a Service}
\fancyfoot[C]{\thepage}
\renewcommand{\headrulewidth}{0.4pt}

% Line spacing
\onehalfspacing

% Bibliography file (create a corresponding .bib file)
\addbibresource{references.bib}

% Glossary entries
\makeglossaries
\newglossaryentry{llm}{name={LLM}, description={Large Language Model}}
\newglossaryentry{cpm}{name={CPM}, description={Cost Per Mille (cost per thousand impressions)}}
\newglossaryentry{ctr}{name={CTR}, description={Click-Through Rate}}
\newglossaryentry{rag}{name={RAG}, description={Retrieval-Augmented Generation}}
\newglossaryentry{ftc}{name={FTC}, description={Federal Trade Commission}}

% ============================================================
% TITLE PAGE
% ============================================================
\begin{document}

\begin{titlepage}
    \centering
    \vspace*{2cm}

    {\Huge\bfseries Ad LLM as a Service}\\[0.5cm]
    {\LARGE An LLM-Powered Conversational Advertising Platform}\\[2cm]

    {\Large\itshape A Project Report}\\[1cm]

    {\large
        \textbf{Author Name}\\
        Student ID: XXXXXXXX\\[0.5cm]
        \textbf{Teammate Name}\\
        Student ID: YYYYYYYY\\[1.5cm]
        School / Department Name\\
        University Name\\[1cm]
        Module Code --- Module Title\\[1.5cm]
        \today
    }

    \vfill
    \textit{Submitted in partial fulfilment of the requirements for [Degree/Module].}
\end{titlepage}

% ============================================================
% FRONT MATTER
% ============================================================
\pagenumbering{roman}

% Abstract
\begin{abstract}
\noindent
% TODO: Write 150–250 word abstract covering:
% - Problem: Traditional display ads are poorly suited to conversational AI interfaces.
% - Approach: An LLM-powered ad platform that extracts user intent from conversations,
%   semantically matches relevant ads, and instructs the LLM to naturally weave sponsored
%   product recommendations into helpful responses.
% - Key components: context extraction, vector-based ad matching, prompt-engineered
%   response synthesis with mandatory disclosure, and engagement analytics.
% - Results/contribution: A working prototype demonstrating higher engagement potential
%   through conversational ad delivery versus traditional impressions.
\end{abstract}

\newpage
\tableofcontents
\newpage
\listoffigures
\listoftables
\listofalgorithms
\newpage

% ============================================================
% MAIN BODY
% ============================================================
\pagenumbering{arabic}

% ============================================================
\section{Introduction}
\label{sec:introduction}
% ============================================================

\subsection{Background and Motivation}
\label{subsec:background}
% - The rise of conversational AI interfaces (ChatGPT, Claude, etc.)
% - The shift in user behaviour from search-engine queries to conversational queries
% - The advertising industry's challenge: traditional banner/display ads do not
%   translate to chat-based interfaces
% - Opportunity: monetise conversational AI through contextually relevant,
%   naturally integrated product recommendations

\subsection{Project Objectives}
\label{subsec:objectives}
% - Design and build an end-to-end LLM-powered conversational ad platform
% - Deliver inline sponsored recommendations that feel like genuine suggestions
% - Ensure FTC-compliant disclosure of sponsored content
% - Implement semantic ad matching using vector embeddings and business-rule re-ranking
% - Track user engagement beyond clicks --- measure conversational follow-ups
% - Produce a functional MVP demonstrating the core value proposition

\subsection{Scope and Constraints}
\label{subsec:scope}
% - MVP scope: single-service architecture, pre-loaded ad inventory, no advertiser portal
% - Constraints: reliance on third-party LLM APIs, latency budgets (<3s total),
%   cost considerations (embedding + two LLM calls per request)

\subsection{Report Organisation}
\label{subsec:report-org}
% Briefly outline what each subsequent section covers.


% ============================================================
\section{Literature Review}
\label{sec:literature-review}
% ============================================================

\subsection{Evolution of Digital Advertising}
\label{subsec:evolution-ads}
% - From banner ads to programmatic advertising
% - Search advertising (Google Ads) and intent-based targeting
% - Social media advertising and behavioural targeting
% - The move towards native advertising and content-integrated formats

\subsection{Large Language Models and Conversational AI}
\label{subsec:llm-background}
% - Transformer architecture and the emergence of LLMs (GPT, Claude, etc.)
% - Capabilities relevant to this project: instruction following, structured output
%   generation, contextual reasoning, tone adaptation
% - Lightweight vs.\ capable models --- the spectrum from Haiku-tier to Sonnet-tier
%   and when each is appropriate

\subsection{Retrieval-Augmented Generation (RAG)}
\label{subsec:rag}
% - RAG paradigm: retrieve relevant context, augment the LLM prompt
% - Vector embeddings and semantic similarity search
% - Embedding models (OpenAI text-embedding-3-small, Cohere Embed, etc.)
% - Vector databases and extensions (pgvector, Pinecone, Weaviate)
% - Relevance to this project: ad matching as a specialised RAG pipeline

\subsection{Advertising in Conversational Interfaces}
\label{subsec:conversational-ads}
% - Existing approaches: Bing Chat ads, Google SGE sponsored results
% - Academic work on native advertising and user perception
% - Ethical considerations: transparency, user trust, ad receptivity
% - Regulatory frameworks: FTC disclosure requirements for sponsored content

\subsection{Prompt Engineering for Controlled Generation}
\label{subsec:prompt-engineering}
% - Techniques for reliable structured output from LLMs
% - Instruction-following and system prompts
% - Balancing helpfulness with ad integration naturalness
% - LLM-as-judge paradigm for automated quality evaluation

\subsection{User Engagement Metrics in Conversational Systems}
\label{subsec:engagement-metrics}
% - Traditional ad metrics: impressions, CTR, conversion rate
% - Novel conversational metrics: follow-up engagement, dwell time, attention verification
% - Why conversational engagement is a stronger signal than click-through
% - Related work on measuring user attention in dialogue systems

\subsection{Summary and Research Gap}
\label{subsec:lit-summary}
% - Identify the gap: no open, well-documented architecture for LLM-native
%   conversational advertising with semantic matching, natural integration,
%   and engagement-based measurement
% - Position your project as addressing this gap


% ============================================================
\section{Problem Statement}
\label{sec:problem-statement}
% ============================================================

\subsection{Limitations of Traditional Ad Formats in Conversational AI}
\label{subsec:traditional-ad-limitations}
% - Banner and sidebar ads are absent in chat interfaces
% - Pre-roll / interstitial formats break conversational flow
% - Search-ad models rely on short keyword queries, not multi-turn dialogues

\subsection{Technical Challenges}
\label{subsec:technical-challenges}

\subsubsection{Real-Time Intent Extraction from Multi-Turn Conversations}
% - User intent is distributed across multiple turns, not a single query
% - Need to classify purchase signals, ad receptivity, and sensitivity in real time

\subsubsection{Semantic Ad Matching at Low Latency}
% - Matching ads to conversational context is more nuanced than keyword matching
% - Balancing relevance, bid price, and frequency caps in a scoring function
% - Latency constraint: the full pipeline must complete in <3 seconds

\subsubsection{Natural Ad Integration Without Degrading User Experience}
% - The LLM must weave product mentions naturally, not as appended blocks
% - Mandatory disclosure ([Sponsored]) must be consistent yet non-intrusive
% - Irrelevant ads must be gracefully rejected by the LLM

\subsubsection{Engagement Measurement Beyond Clicks}
% - Clicks are sparse in conversational settings
% - Need to detect whether a user's follow-up message engages with the shown ad
% - This requires a classification capability (keyword matching → LLM classifier)

\subsection{Ethical and Regulatory Considerations}
\label{subsec:ethical-considerations}
% - FTC compliance: clear and conspicuous disclosure of sponsored content
% - User trust: ads in a conversational context carry higher trust expectations
% - Sensitivity gating: suppressing ads during health crises, emotional distress,
%   financial distress, or grief
% - Data privacy: session data handling and retention policies

\subsection{Formal Problem Definition}
\label{subsec:formal-problem}
% Optionally, a concise formal statement:
% Given a user message m_t and conversation history H = {m_1, ..., m_{t-1}},
% and an ad inventory A, produce a response r_t that:
% (1) helpfully answers the user's query,
% (2) naturally integrates a relevant sponsored product a* ∈ A if contextually
%     appropriate, with mandatory disclosure, and
% (3) logs engagement signals for performance measurement.


% ============================================================
\section{Solution Design}
\label{sec:solution-design}
% ============================================================

\subsection{System Architecture Overview}
\label{subsec:architecture-overview}
% - High-level architecture diagram (TikZ or included figure)
% - Data flow: User Message → Orchestrator → Context Extraction → Ad Matching
%   → Response Synthesis → Frontend
% - Tech stack summary table

\begin{figure}[H]
    \centering
    % TODO: Replace with your architecture diagram (TikZ or \includegraphics)
    % \includegraphics[width=\textwidth]{figures/architecture.png}
    \fbox{\parbox{0.9\textwidth}{\centering\vspace{3cm}
        \textit{[System Architecture Diagram --- see Section~\ref{subsec:architecture-overview}]}
    \vspace{3cm}}}
    \caption{High-level system architecture of the Ad LLM platform.}
    \label{fig:architecture}
\end{figure}

\begin{table}[H]
    \centering
    \caption{Technology stack summary.}
    \label{tab:tech-stack}
    \begin{tabular}{@{}ll@{}}
        \toprule
        \textbf{Component}          & \textbf{Technology}                   \\ \midrule
        Frontend                    & React                                 \\
        Orchestrator API            & FastAPI (Python)                      \\
        Session Store               & Redis                                 \\
        Database                    & PostgreSQL + pgvector                 \\
        Embedding Model             & text-embedding-3-small (OpenAI)       \\
        Context Extraction LLM      & Claude Haiku (lightweight)            \\
        Response Generation LLM     & Claude Sonnet (capable)               \\
        Event Logging               & PostgreSQL (JSONB payload)            \\
        \bottomrule
    \end{tabular}
\end{table}

\subsection{Component Design}
\label{subsec:component-design}

\subsubsection{Orchestrator API (FastAPI)}
\label{subsubsec:orchestrator}
% - Central routing service: request validation, auth, rate limiting
% - Orchestration flow for POST /chat (the seven-step pipeline)
% - Graceful degradation: if any sub-component fails, return a normal response
% - Latency optimisation: parallel pre-fetching, warm-start vector search
% - Endpoint specifications: /chat, /track/click, /analytics/summary, /admin/ads

\subsubsection{Context Extraction Module}
\label{subsubsec:context-extraction}
% - LLM-based (lightweight model) extraction of structured context from conversation
% - Output schema: intent, topics, entities, sentiment, purchase_signal, ad_receptivity
% - Prompt design rationale and iteration process
% - Safety/sensitivity classifier: ad_receptivity = "none" for sensitive conversations
% - Test set: 30+ conversation snippets with expected outputs

\subsubsection{Session Management (Redis)}
\label{subsubsec:session-management}
% - Session schema: turns (rolling window of 10), accumulated topics,
%   purchase signal history, ads shown (frequency cap), turn count
% - TTL policy: 30-minute inactivity expiry
% - Role in the pipeline: provides conversation history and frequency-cap data

\subsubsection{Ad Inventory Store (PostgreSQL + pgvector)}
\label{subsubsec:ad-inventory}
% - Database schema: ads table (with vector(1536) column), events table, advertisers table
% - pgvector justification over dedicated vector DB (scale, simplicity)
% - IVFFlat index configuration
% - Ad data model: product metadata, targeting fields, bid/budget, embedding

\subsubsection{Embedding and Semantic Matching Pipeline}
\label{subsubsec:embedding-matching}
% - Ad embedding at ingestion time: embed_ad(text) → vector
% - Query embedding at request time: embed_query(context) → vector
% - What gets embedded on each side and why
% - Vector similarity search: cosine distance, top-20 candidates

\subsubsection{Re-Ranking and Business Logic}
\label{subsubsec:reranking}
% - Post-retrieval filtering: active ads, positive budget, frequency cap
% - Composite scoring formula:
%   score = similarity × 0.5 + bid_cpm_normalised × 0.3 + purchase_signal × 0.2
% - Rationale for weight selection
% - Return top 3 candidates to orchestrator

\begin{algorithm}[H]
    \caption{Ad Candidate Re-Ranking}
    \label{alg:reranking}
    \begin{algorithmic}[1]
        \Require Query embedding $\mathbf{q}$, context $C$, session ID $s$
        \Ensure Ranked list of top-$k$ ad candidates
        \State $\mathit{raw} \gets \text{VectorSearch}(\mathbf{q}, \text{limit}=20)$
            \Comment{pgvector cosine similarity}
        \State $\mathit{shown} \gets \text{Redis.get}(s).\text{ads\_shown}$
        \State $\mathit{candidates} \gets \{a \in \mathit{raw} \mid a.\text{active} \land a.\text{budget} > 0
            \land a.\text{id} \notin \mathit{shown}\}$
        \State $b_{\max} \gets \max_{a \in \mathit{candidates}} a.\text{bid\_cpm}$
        \ForAll{$a \in \mathit{candidates}$}
            \State $a.\text{score} \gets 0.5 \cdot a.\text{sim} + 0.3 \cdot \frac{a.\text{bid\_cpm}}{b_{\max}}
                + 0.2 \cdot C.\text{purchase\_signal}$
        \EndFor
        \State \Return $\text{TopK}(\mathit{candidates}, k=3)$
    \end{algorithmic}
\end{algorithm}

\subsubsection{Response Synthesis with Ad Injection}
\label{subsubsec:response-synthesis}
% - Main LLM call (Sonnet-tier) with carefully constructed prompt
% - Prompt structure: system instruction, <sponsored_context> block, conversation history
% - Design principles: naturalness, mandatory [Sponsored] disclosure, parseable
%   markdown links, graceful rejection of irrelevant ads
% - Output schema: response_text, ad_included, ad_id_used
% - Prompt engineering challenges and iteration methodology

\subsubsection{Click Tracking and Event Logging}
\label{subsubsec:tracking}
% - Tracking link generation: wrap CTA URLs through /track/click redirect
% - Event model: impression, click, engagement, conversion
% - Click redirect endpoint: log event → 302 redirect to advertiser URL
% - Impression logging with CPM budget deduction
% - Engagement detection: keyword/entity matching → LLM classifier

\subsubsection{Analytics Dashboard}
\label{subsubsec:analytics}
% - /analytics/summary endpoint: impressions, CTR, engagements, revenue, top ads
% - Internal React dashboard consuming the analytics API
% - Key metrics and their computation

\subsection{Interface Contracts and Workstream Division}
\label{subsec:interface-contracts}
% - Two-person team: Backend/Infrastructure vs.\ AI/LLM
% - Agreed function signatures: extract_context, embed_query, embed_ad,
%   get_candidate_ads, generate_response
% - Communication pattern: backend calls AI modules; AI calls backend ad retrieval

\begin{table}[H]
    \centering
    \caption{Interface contracts between Backend and AI workstreams.}
    \label{tab:interface-contracts}
    \begin{tabular}{@{}p{3.5cm}p{4cm}p{5cm}@{}}
        \toprule
        \textbf{Function}           & \textbf{Caller → Callee}  & \textbf{Description}              \\ \midrule
        \texttt{extract\_context}   & Backend → AI              & Extract structured context from
                                                                  user message and history           \\
        \texttt{embed\_query}       & Backend → AI              & Generate query embedding from
                                                                  extracted context                  \\
        \texttt{embed\_ad}          & Backend → AI              & Generate ad embedding at
                                                                  ingestion time                     \\
        \texttt{get\_candidate\_ads}& AI → Backend              & Retrieve and rank ad candidates
                                                                  via vector search                  \\
        \texttt{generate\_response} & Backend → AI              & Synthesise response with optional
                                                                  ad injection                       \\
        \bottomrule
    \end{tabular}
\end{table}

\subsection{Quality Assurance and Evaluation Framework}
\label{subsec:quality-assurance}
% - LLM-as-judge evaluator: helpfulness, naturalness, disclosure accuracy, ad relevance
% - Evaluation sets for context extraction (30+ scenarios) and response synthesis (20–30 scenarios)
% - Engagement detection accuracy testing
% - End-to-end integration testing

\subsection{Error Handling and Graceful Degradation}
\label{subsec:error-handling}
% - Failure modes and fallback strategies for each component
% - Timeout configuration: 10s for context extraction, 30s for response generation
% - Monitoring and alerting: error rate, latency, ad fill rate


% ============================================================
\section{Illustrative Examples}
\label{sec:illustrative-examples}
% ============================================================

\subsection{Example 1: Product Research Conversation}
\label{subsec:example-product-research}
% - Scenario: User asks "What are the best running shoes for marathon training?"
% - Walk through the full pipeline:
%   1. Context extraction output (intent: product_research, purchase_signal: 0.85,
%      ad_receptivity: high)
%   2. Query embedding and vector search results
%   3. Re-ranking with business rules → selected ad (e.g., Nike Pegasus 41)
%   4. Prompt assembled with <sponsored_context>
%   5. Generated response with natural [Sponsored] mention and tracking link
%   6. Impression logged, budget deducted
% - Show the actual JSON payloads at each step

\subsection{Example 2: Sensitive Conversation (Ad Suppression)}
\label{subsec:example-sensitive}
% - Scenario: User is discussing financial distress or emotional difficulty
% - Context extraction returns ad_receptivity: "none"
% - Pipeline skips ad matching entirely
% - Response generated without any <sponsored_context> block
% - Demonstrates the ethical safeguard

\subsection{Example 3: Irrelevant Ad Rejection}
\label{subsec:example-irrelevant}
% - Scenario: User asks about cooking; the only available ads are for running shoes
% - Context extraction returns relevant topics (cooking, recipes)
% - Vector search returns low-similarity running shoe ads
% - Even if passed to the LLM, the prompt instructs it to ignore irrelevant ads
% - Generated response is purely helpful with no product mention

\subsection{Example 4: Engagement Follow-Up}
\label{subsec:example-engagement}
% - Scenario: After seeing an ad for Nike Pegasus 41, user asks
%   "Tell me more about the Pegasus 41"
% - Engagement classifier detects follow-up about the advertised product
% - Engagement event logged (stronger signal than a click)
% - Subsequent response provides more detail, potentially with a second ad mention

\subsection{Example 5: Multi-Turn Context Accumulation}
\label{subsec:example-multi-turn}
% - Scenario: Conversation starts casual, then shifts towards product research
%   over several turns
% - Show how accumulated_topics and purchase_signal_history evolve in Redis
% - Show how ad_receptivity transitions from "low" → "medium" → "high"
% - Demonstrate the session state driving dynamic ad targeting


% ============================================================
\section{Conclusions}
\label{sec:conclusions}
% ============================================================

\subsection{Summary of Contributions}
\label{subsec:summary-contributions}
% - End-to-end architecture for LLM-native conversational advertising
% - Semantic ad matching via RAG-style vector retrieval + business-rule re-ranking
% - Prompt-engineered natural ad integration with mandatory FTC-compliant disclosure
% - Novel engagement metric: conversational follow-up as a stronger signal than CTR
% - Ethical safeguards: sensitivity-aware ad gating

\subsection{Workstream Contributions}
\label{subsec:workstream-contributions}
% - Backend/Infrastructure: orchestrator, session management, database schema,
%   ad matching pipeline, click tracking, analytics API, frontend integration
% - AI/LLM: context extraction, embedding pipeline, response synthesis,
%   engagement classifier, LLM-as-judge evaluator

\subsection{Limitations}
\label{subsec:limitations}
% - Dependence on third-party LLM APIs (latency, cost, availability)
% - Heuristic re-ranking weights (not learned from data)
% - Small ad inventory for testing; scalability not yet validated at 100k+ ads
% - No real advertiser feedback or A/B testing results
% - Engagement classifier starts with keyword matching (limited accuracy)

\subsection{Future Work}
\label{subsec:future-work}
% - Attention verification model (dwell time, read-through estimation)
% - Predicted CTR model trained on click/engagement data
% - Advertiser self-serve portal with campaign management
% - A/B testing infrastructure for ad insertion strategies
% - Tone adaptation: match ad copy to conversational register
% - Brand safety and content filtering (topic classifier + blocklist)
% - Fine-tuning context extraction on domain-specific data

\subsection{Reflections}
\label{subsec:reflections}
% - Lessons learned from the two-workstream collaboration
% - Prompt engineering iteration process and key insights
% - Balancing user experience with monetisation goals
% - Ethical considerations encountered during development


% ============================================================
% REFERENCES
% ============================================================
\newpage
\printbibliography[heading=bibintoc, title={References}]

% Placeholder: create a references.bib file with entries such as:
% @article{vaswani2017attention,
%   title={Attention is all you need},
%   author={Vaswani, Ashish and others},
%   journal={Advances in Neural Information Processing Systems},
%   year={2017}
% }
% @misc{openai2023gpt4,
%   title={GPT-4 Technical Report},
%   author={{OpenAI}},
%   year={2023}
% }
% @article{lewis2020rag,
%   title={Retrieval-augmented generation for knowledge-intensive NLP tasks},
%   author={Lewis, Patrick and others},
%   journal={Advances in Neural Information Processing Systems},
%   year={2020}
% }
% @misc{ftc2023endorsements,
%   title={FTC Endorsement Guides},
%   author={{Federal Trade Commission}},
%   year={2023},
%   url={https://www.ftc.gov/legal-library/browse/federal-register-notices/16-cfr-part-255-guides-concerning-use-endorsements-testimonials-advertising}
% }


% ============================================================
% APPENDICES
% ============================================================
\newpage
\begin{appendices}

\section{Database Schema}
\label{app:schema}
% - Full CREATE TABLE statements for ads, events, advertisers
% - Index definitions
% - Include the SQL from the backend workstream document

\section{API Endpoint Specifications}
\label{app:api-endpoints}
% - Full request/response schemas for all endpoints
% - POST /chat, GET /track/click, GET /analytics/summary, POST /admin/ads

\section{Prompt Templates}
\label{app:prompts}
% - Context extraction prompt (full text)
% - Response synthesis system prompt (full text)
% - LLM-as-judge evaluation prompt (full text)

\section{Context Extraction Test Cases}
\label{app:test-cases-extraction}
% - Table of 30+ test conversation snippets with expected extraction outputs

\section{Response Synthesis Evaluation Set}
\label{app:test-cases-synthesis}
% - Table of 20–30 scenarios with expected quality scores

\section{Session State Schema}
\label{app:session-schema}
% - Full Redis session object schema with field descriptions

\section{Project Timeline}
\label{app:timeline}
% - Gantt chart or week-by-week table for both workstreams

\begin{table}[H]
    \centering
    \caption{Project timeline by workstream.}
    \label{tab:timeline}
    \begin{tabular}{@{}clll@{}}
        \toprule
        \textbf{Week} & \textbf{Backend Workstream}    & \textbf{AI/LLM Workstream}    & \textbf{Milestone}    \\ \midrule
        1   & Orchestrator skeleton, Redis,      & Context extraction prompt      & Infrastructure        \\
            & Postgres + pgvector schema         & design + module                & foundation            \\[4pt]
        2   & Ad CRUD API, bulk ingestion,       & Embedding functions,           & Ad matching           \\
            & re-ranking logic                   & retrieval quality testing       & pipeline ready        \\[4pt]
        3   & Click tracking, event logging,     & Response synthesis prompt       & Core pipeline         \\
            & tracking link generator            & engineering + module           & complete              \\[4pt]
        4   & Analytics API, frontend            & Eval set, quality tuning,      & End-to-end            \\
            & integration, end-to-end wiring     & engagement classifier          & integration           \\[4pt]
        5   & Error handling, monitoring,        & LLM-as-judge evaluator,        & Demo-ready            \\
            & hardening                          & attention model start          & MVP                   \\
        \bottomrule
    \end{tabular}
\end{table}

\end{appendices}

% ============================================================
% GLOSSARY
% ============================================================
\newpage
\printglossaries

\end{document}